\chapter{Background}\label{chap:background}

This chapter provides the terminology and prior work needed to follow the methods (Chapter~\ref{chap:methods}) and experiments (Chapter~\ref{chap:experiments}): the WTI market and its liquidity measures, the NLP tools used to turn news into features, the Temporal Fusion Transformer, and the strands of literature that motivate the design choices of this thesis.

\section{Terminology}\label{sec:terminology}

This section fixes the vocabulary used throughout; each concept is developed at its point of use in Chapters~\ref{chap:methods} and~\ref{chap:experiments}.

\textbf{WTI crude oil futures and the EIA inventory cycle.} West Texas Intermediate (WTI) is the U.S.\ benchmark grade of crude oil; this thesis studies its front-month futures contract, the most actively traded WTI instrument. WTI is chosen for its deep liquidity, its sensitivity to a well-defined event structure (geopolitics, OPEC policy, macroeconomic releases), and the weekly cadence of the U.S.\ Energy Information Administration (EIA) Weekly Petroleum Status Report, which publishes national crude inventory levels every Wednesday at 10:30 a.m.\ Eastern Time and is the most regular scheduled event in the oil calendar.

\textbf{Liquidity measures.} Market liquidity is the ease of trading an asset without moving its price; this thesis quantifies it hourly with three complementary measures. Trading volume, used in log form, captures raw activity. The Amihud illiquidity ratio \citep{amihud2002}, the absolute return per unit of volume, captures price impact, how far the price moves for a given amount of trading. The Parkinson range \citep{parkinson1980}, the log high-low range within an hour, is a high-frequency estimator of volatility. Together they describe activity, price impact, and volatility.

\textbf{NLP for finance.} Turning news text into numerical features draws on three natural-language-processing ideas. Three-class sentiment classification labels a text as positive, neutral, or negative, the standard output of financial sentiment models. Structured extraction instead asks a model to fill a fixed schema of fields (for example a sentiment score, an event type, and a list of entities), yielding richer, multi-dimensional features. Prompt engineering, and its stricter form of tool-use schema enforcement, is how a large language model is instructed to produce that structured output reliably.

\textbf{Temporal Fusion Transformer.} The Temporal Fusion Transformer (TFT) \citep{lim2021} is an attention-based neural architecture for multi-horizon time-series forecasting. Two of its components give this thesis its interpretability: the Variable Selection Network, which assigns each input feature a learned importance weight, and the interpretable multi-head attention layer, which exposes how much the model attends to each past time step when forecasting. These let the model's reasoning be read off directly rather than treated as a black box.

\section{Related work}\label{sec:related-work}

The design decisions of this thesis are grounded in prior work rather than derived in isolation. This section states the anchors; each is picked up at the point of use in Chapters~\ref{chap:methods} and~\ref{chap:experiments}.

\textbf{News sentiment and market activity.} The premise that news sentiment moves trading activity, and not only prices, is established in the finance literature: media pessimism has been shown to predict downward price pressure and elevated trading volume \citep{tetlock2007}. This grounds the use of news sentiment as a predictor of hourly liquidity in Phase~1 (Chapter~\ref{chap:experiments}, Section~\ref{sec:phase1}).

\textbf{News at intraday resolution.} The market response to news is concentrated at high frequency rather than spread evenly across the day: scheduled macroeconomic announcements account for most of the intraday and day-of-week volatility patterns in interest-rate and foreign-exchange futures, with the bulk of the adjustment occurring within minutes of release and volatility remaining elevated for several hours afterward \citep{ederington1993}. Most news-sentiment studies in finance nonetheless operate at daily resolution, aggregating away exactly the intraday structure this thesis sets out to measure. Working at hourly resolution (Chapter~\ref{chap:methods}) is the central methodological choice that follows.

\textbf{News in oil markets.} The price response of oil to news is muted at coarse resolution: daily energy-price changes show no compelling response to daily U.S.\ macroeconomic news releases, consistent with energy prices being predetermined with respect to macro aggregates at the daily horizon \citep{kilian_vega2011}. A null response at daily frequency does not preclude one at finer resolution, and it motivates this thesis's twin choices of hourly data and a liquidity target (trading volume) rather than price.

\textbf{Financial-text sentiment models.} Domain-adapted transformer models outperform general-purpose sentiment tools on financial language; FinBERT \citep{araci2019} is the de facto baseline for financial sentiment and is adopted as the Phase~1 sentiment extractor (Section~\ref{sec:finbert}). Its coarse three-class output is also the limitation that motivates the richer LLM-based extraction of Phase~2.

\textbf{Liquidity measurement.} Hourly liquidity is quantified with standard market-microstructure measures: trading volume, the Parkinson high-low range \citep{parkinson1980}, and the Amihud illiquidity ratio \citep{amihud2002}, which captures price impact per unit of traded volume.

\textbf{Oil-price shock decomposition.} Oil-market movements are not a single undifferentiated shock but a composition of separable drivers: the structural decomposition of the real oil price into supply, aggregate-demand, and precautionary (risk) demand components \citep{kilian2009} is the canonical framework. This grounds the Phase~2 channel decomposition (Sections~\ref{sec:channel-decomposition} and~\ref{sec:llm-extraction}), which extracts a \texttt{supply\_impact}, \texttt{demand\_impact}, and \texttt{risk\_premium} from each article rather than collapsing them onto a single sentiment axis, and it clarifies why supply-risk geopolitical news is price-bullish rather than bearish.

\textbf{Interpretable multi-horizon forecasting.} The Temporal Fusion Transformer \citep{lim2021} is an attention-based architecture for multi-horizon time-series forecasting that couples predictive accuracy with interpretability through a Variable Selection Network and interpretable multi-head attention. These two mechanisms are what let Phase~2 (Sections~\ref{sec:tft-methods}, \ref{sec:tft-v1}, and~\ref{sec:tft-v2}) read feature importance and temporal attention off the model instead of treating it as a black box, and the thesis uses them directly to corroborate the lag structure of RQ1.

\textbf{LLM-based extraction and cross-model validation.} Large language models can extract structured, multi-dimensional features from financial text, and they can stand in for costly human annotation, matching or exceeding crowd workers on text-labelling tasks \citep{gilardi2023}, and can serve as evaluators of other models' outputs \citep{zheng2023}. Phase~2 uses both ideas: Claude Haiku extracts the structured schema of Section~\ref{sec:schema}, producing the richer, price-oriented, channel-decomposed features that FinBERT's three-class tone label cannot, and the inter-model calibration of Section~\ref{sec:calibration-methods} (Section~\ref{sec:calibration}) validates those features by cross-model agreement against a second model.
